\begin{figure}[!t]
\centering
\begin{minipage}[t]{.47\linewidth}
\vspace{0pt}\centering
\textbf{(a) Semantic computation graph}\\[9pt]
% A toy speed-control subgraph, paired with speed_control.py.
\begin{tikzpicture}[
    x=1cm,y=1cm,>=Stealth,
    var/.style={draw,line width=.65pt,inner sep=3pt,font=\ttfamily\fontsize{8}{9}\selectfont},
    obs/.style={var,draw=black!60,fill=black!6},
    latent/.style={var,draw=blue!65!black,fill=blue!9},
    action/.style={var,draw=green!40!black,fill=green!12},
    weight/.style={var,draw=orange!75!black,fill=orange!16},
    op/.style={draw=black!65,dashed,rounded corners=4pt,
        line width=.65pt,fill=white,inner sep=4pt,font=\fontsize{8}{9}\selectfont,align=center},
    flow/.style={->,line width=.65pt},
    key/.style={font=\fontsize{6.5}{8}\selectfont,inner sep=2pt}
]
\node[weight] (target) at (0,0) {target\_speed};
\node[obs] (speed) at (2,0) {speed};
\node[op] (subtract) at (1,-.8) {subtract};
\node[latent] (error) at (1,-1.55) {speed\_error};
\node[weight] (gain) at (-1.6,-1.55) {gain};
\node[op] (control) at (-.3,-2.5) {ReLU $\to$ multiply\\$\to$ clamp to $[0,1]$};
\node[action] (throttle) at (-.3,-3.45) {throttle};
\draw[flow] (target) -- node[left,font=\scriptsize] {} (subtract);
\draw[flow] (speed) -- node[right,font=\scriptsize] {} (subtract);
\draw[flow] (subtract) -- (error);
\draw[flow] (error) -- (control);
\draw[flow] (gain) -- (control);
\draw[flow] (control) -- (throttle);
\matrix[column sep=3pt] at (-.3,-4.15) {
    \node[obs,key] {Observation}; &
    \node[latent,key] {Latent}; &
    \node[action,key] {Action}; &
    \node[weight,key] {Weight (trainable)}; \\
};
\end{tikzpicture}

\end{minipage}\hfill
\begin{minipage}[t]{.51\linewidth}
\vspace{0pt}\centering
\textbf{(b) Corresponding PyTorch code snippet}\\[5pt]
\lstinputlisting[language=Python,firstline=4,
    basicstyle=\ttfamily\fontsize{8.5}{10}\selectfont,
    emph={target_speed,gain},emphstyle=\color{orange!65!black},
    emph={[2]speed_error},emphstyle={[2]\color{blue!65!black}},
    emph={[3]throttle},emphstyle={[3]\color{green!40!black}}
]{figures/diagrams/speed_control.py}
\end{minipage}
\caption{
A sample of speed-control subgraph and its correspondingPyTorch implementation.
Solid nodes distinguish observations (gray), semantic latents (blue), action outputs (green), and trainable weights (orange); dashed nodes denote operation nodes (some collapsed due to space constraints).
The returned \texttt{speed\_error} exposes the latent value for further rollout analysis.
}
\label{fig:topology}
\end{figure}

\section{Semantic policy structure}
\label{sec:structure}
\label{sec:policy-program}
\label{sec:running-example}
\label{sec:semantic-interface}

Following previous work \cite{zhu2025sample}, a structured policy computes actions from observations through a semantically meaningful latent.
Formally, each structured policy has a symbolic component $\llbracket\alpha\rrbracket = \langle (V, P), E \rangle$: a bipartite graph whose variable nodes $V$ represent intermediate values, operator nodes $P$ compute transformations between them, and edges $E$ connect variables and operators.
Each structured policy also has a neural component $\Theta$ containing all learnable latent values (a subset of $V$), which determine the exact continuous output.
We implement each structured policy as a task-specific PyTorch model \cite{paszke2019pytorch}, representing the graph topology in code and fitting its parameters through the autograd function.
We denote the policy with structure $\llbracket \alpha \rrbracket$ and trained weights $\Theta$ by $\pi_{\langle \llbracket \alpha \rrbracket, \Theta \rangle}$, using $\llbracket \alpha \rrbracket$ for both the graph topology and its PyTorch implementation.

Figure~\ref{fig:topology} shows a minimal speed-control example alongside its underlying implementation.
The topology defines a proportional controller for maintaining constant speed, while imitation learning fits parameters such as \texttt{target\_speed} and \texttt{gain} to the demonstration's speed and responsiveness.
This highlights the separation of symbolic and neural components: the symbolic structure is task-specific, specifying the intermediate steps that generate actions; continuous parameters are demonstration-specific, determining how the demonstration completes the task.

Unlike post-hoc explanations of black-box models, latent semantics directly reveal expected policy behavior, while intermediate computations trace how actions are generated.
In this work, we audit all the internal mechanisms that generate actions and use this information to detect structure-demonstration misalignment and propose structural refinement.